% Chapter 2

\chapter{Review of Related Literature} % Main chapter title

\label{Chapter2}

%----------------------------------------------------------------------------------------

\section{Introduction}

This chapter reviews the knowledge on which the proposed platform builds. It starts with conversational agents in
customer service and their effect on customers and staff, continues with large language models used as agents that call
tools, and then discusses the main risk of such agents --- invented facts --- and how grounding reduces it. It then
reviews the state of natural language processing for Kyrgyz, the conditions of the Kyrgyz market, software-as-a-service
architecture for many small tenants, and existing products. The chapter ends with a summary and the research gap.

\section{Conversational Agents in Customer Service}

Chatbots have evolved from pattern-matching programs to systems based on machine learning and, most recently, on large
language models \parencite{adamopoulou2020chatbots}. Customer service on social media was an early target:
\textcite{xu2017new} built a chatbot for customer requests on Twitter, found that over 40\% of the requests were
emotional rather than informational, and reported that their system was about as good as human agents in showing empathy.
\textcite{folstad2017chatbots} argued that chatbots change interaction design itself: the conversation becomes the user
interface, so the quality of the dialogue, not the layout of screens, decides whether the service works.

The business effect of such agents depends on how customers react to them. \textcite{adam2021ai} showed in an online
experiment that anthropomorphic design cues and the foot-in-the-door technique increased users' compliance with a
chatbot's requests in customer service. For a small business this means that the agent's wording --- greeting, asking
for a phone number, confirming a booking --- directly affects whether a conversation ends with a booking.

Evidence from the field is now available for LLM assistants. \textcite{brynjolfsson2025generative} studied the staggered
introduction of a generative AI assistant among customer support agents and found a 15\% average increase in issues
resolved per hour, with the largest gains for novice and lower-skilled workers. The setting differs from the one in this
thesis --- the assistant suggested answers to human agents rather than answering customers directly --- but it shows
that LLMs can carry the knowledge of experienced staff into everyday customer conversations.

\section{Large Language Models as Agents}

\subsection{From Language Models to Agents}

Large language models learn to perform new tasks from instructions and a few examples in the prompt, without retraining
\parencite{brown2020language}. An \emph{agent} built on an LLM additionally perceives an environment, plans and acts in it
\parencite{wang2024survey}. For customer service, the environment is the business: its price list, its staff and their
schedules, its free time slots, its stock and its open requests.

\subsection{Reasoning and Tool Use}

\textcite{yao2023react} showed that interleaving reasoning with actions (the ReAct approach) lets a model look up
information and act on it, reducing hallucination compared with reasoning alone. \textcite{schick2023toolformer} showed
that a model can learn when to call external tools such as calculators and search engines. Commercial model APIs now
expose this capability as structured outputs and function calling: the model returns a machine-readable action, and the
application decides whether to execute it.

\subsection{Evaluating Agents}

Evaluating agents that talk to users and act on a database is harder than evaluating single answers.
\textcite{yao2024taubench} proposed $\tau$-bench, in which an agent serves a simulated user in retail and airline domains
under the domain's policies, and the result is judged by the final state of the database rather than by the wording of
the reply. They report that even strong models complete fewer than half of the tasks and are inconsistent across
repeated trials. \textcite{zheng2023judging} studied the use of a strong LLM as a judge of chatbot answers and found high
agreement with human preferences, but also position, verbosity and self-enhancement biases. Both works support evaluating
a business agent with explicit, checkable outcomes --- was a booking created, was a fact invented --- in addition to
judgments of answer quality.

\section{Hallucination and Grounding}

LLMs may generate content that is fluent but unfaithful to the source or to the world; \textcite{ji2023survey} survey this
\emph{hallucination} problem across tasks and mitigation methods. Retrieval-augmented generation
\parencite{lewis2020retrieval} conditions the model on documents retrieved for the query and improves factuality on
knowledge-intensive tasks. For a single small business the relevant knowledge is small enough to be given to the model
directly: the price list, working hours, address, staff and frequent questions fit into the prompt. Grounding alone,
however, does not prevent the model from \emph{claiming} an action that did not happen, such as confirming a booking
for a time that is already taken. This motivates validating actions in code and not only grounding the text.

\section{Natural Language Processing for Kyrgyz}

\textcite{joshi2020state} classified the world's languages by the amount of labelled and unlabelled data available and
showed that most languages are left behind by NLP progress. In a recent overview, \textcite{alekseev2024kyrgyznlp}
describe Kyrgyz as a language that is only \enquote{scraping by} in resources: apart from dictionaries, most public
datasets and tools appeared only recently, and the authors propose a roadmap for community-driven development.
\textcite{turatali2026kyrgyzllmbench} introduced KyrgyzLLM-Bench, with natively written and manually post-edited
translated tasks, and evaluated 26 open and proprietary LLMs; model rankings transfer from English to Kyrgyz on some tasks
but not others, and several tasks show a substantial gap between English and Kyrgyz.

For customer service, an additional difficulty is code-switching: customers in Kyrgyzstan write in Kyrgyz, in Russian, or
mix both, and may use Latin transliteration. A business agent must therefore detect the language of each message, answer
in the same language and know the Kyrgyz names of the business's services, while keeping prices and facts identical in
both languages.

\section{Messaging Channels and the Kyrgyz Market}

Small and medium-sized businesses produced 47.3\% of the GDP of Kyrgyzstan in the first half of 2025, according to the
National Statistics Committee \parencite{akchabar2025sme}. Internet use reached 88.5\% of the population at the end of
2025, and Instagram's advertising reach was 3.90~million users \parencite{datareportal2026kg}. Businesses use Instagram
as a shop window and WhatsApp and Telegram for conversations; each channel offers a business API with its own rules for
webhooks, message windows and media. An agent for this market therefore has to work across several channels at once and
keep one customer history across them.

\section{Software-as-a-Service for Many Small Businesses}

A product for small businesses has to be sold and operated as a service: one deployment serves many businesses
(tenants), each configured by its owner without the developer's help. \textcite{bezemer2010multi} discuss the
trade-offs of multi-tenant software-as-a-service: lower operating cost and easier updates, at the price of strict data
isolation between tenants and configurability within one code base. For an LLM agent, tenant isolation also applies to
the prompt: the agent of one business must never see another business's data.

\section{Existing Solutions and Their Limitations}

Products available to small businesses in the region fall into three groups.
\begin{itemize}
	\item \textbf{Online booking and CRM systems} for service businesses keep schedules, clients, loyalty programmes and
	reports, and offer a booking page. They manage records well, but answering a customer's message in a messenger
	remains a human task.
	\item \textbf{Omnichannel chat platforms} collect messages from several messengers into one inbox and provide team
	features such as ready-made answer templates, internal comments and block lists. They make human agents faster but
	still require a human to answer.
	\item \textbf{General-purpose chatbots} and chatbot builders can answer automatically, but they either follow fixed
	scripts that break on free-form questions, or use an LLM without validating actions against the business's real
	schedule, which risks invented confirmations.
\end{itemize}
None of these groups combines automatic answers in Kyrgyz and Russian with validated bookings and orders and a
self-service setup for a non-technical owner.

\section{Summary of the Reviewed Literature}

The literature shows that conversational agents can serve customers and that LLM assistants raise productivity in
customer support \parencite{brynjolfsson2025generative}; that agents which call tools are powerful but hard to make
reliable \parencite{yao2023react, yao2024taubench}; that hallucination is a central risk and grounding reduces but does not
remove it \parencite{ji2023survey, lewis2020retrieval}; and that Kyrgyz remains a low-resource language for which
evaluation resources are only now appearing \parencite{alekseev2024kyrgyznlp, turatali2026kyrgyzllmbench}.

\section{Research Gap}

Published work on LLM agents for customer service concentrates on English and on large companies with dedicated support
teams, and evaluates agents mostly in simulation. There is little work on agents that answer customers of
\emph{small} businesses directly, in a \emph{low-resource language} with code-switching, and that are constrained to
act only on validated business data. There is also no public, bilingual Kyrgyz--Russian evaluation set for such a
customer-service agent. This thesis addresses the gap by designing, implementing and evaluating such a platform and by
publishing its evaluation set.
